% -------------------------
%
% AGRADECIMIENTOS (recomendable máx. 1 pág.)
%
% -------------------------
\phantomsection % Necesario con hyperref
\chapter*{Agradecimientos} % Opción con * para que no aparezca numerada
\addcontentsline{toc}{chapter}{Agradecimientos} % Añade al TOC.

Aunque no es un apartado obligatorio, dedicar unas líneas a expresar agradecimiento es una buena forma de reconocer el apoyo recibido a lo largo del proceso. Este espacio permite dar las gracias a todas aquellas personas que, de una u otra manera, han contribuido a que este trabajo haya podido desarrollarse y finalizar con éxito. Es habitual mencionar aquí a tutores, profesores, familiares, compañeros o amistades que hayan acompañado durante el camino.

Los agradecimientos pueden tener un tono personal e incluir alguna anécdota si se desea, pero se recomienda que no excedan una página de extensión.



\begin{flushright}
  \vspace{1.5cm}  % con punto, no coma
  \textit{Nombre del autor}\\
  Huelva, 2025
\end{flushright}